\corrigesChapitre{20}{Signaux et information}
\par\noindent\begin{minipage}{\linewidth}
\begin{corrige}{1}{Les rôles}
\begin{questions}
\item L’avertisseur émet ; l’oreille de l’utilisateur reçoit le signal.
\item Une demande d’attention ou l’annonce d’une présence, selon la convention.
\end{questions}
\end{corrige}
\end{minipage}
\par\noindent\begin{minipage}{\linewidth}
\begin{corrige}{2}{Radio et son}
\begin{questions}
\item Non : c’est un rayonnement électromagnétique.
\item Oui, contrairement au son qui nécessite un milieu matériel.
\end{questions}
\end{corrige}
\end{minipage}
\par\noindent\begin{minipage}{\linewidth}
\begin{corrige}{3}{Deux états}
\begin{questions}
\item Deux valeurs ou catégories distinguables représentent par exemple 0 et 1.
\item Non : le sens dépend de la convention annoncée.
\end{questions}
\end{corrige}
\end{minipage}
\par\noindent\begin{minipage}{\linewidth}
\begin{corrige}{4}{Le capteur}
\begin{questions}
\item Relier une grandeur physique à un signal exploitable.
\item Pour établir ou vérifier la correspondance entre la grandeur et la sortie.
\end{questions}
\end{corrige}
\end{minipage}
\par\noindent\begin{minipage}{\linewidth}
\begin{corrige}{5}{Une perturbation}
\begin{questions}
\item Un obstacle, une mauvaise orientation ou un éclairage ambiant gênant.
\item Un obstacle, une source de bruit ou une distance trop grande.
\end{questions}
\end{corrige}
\end{minipage}
\par\noindent\begin{minipage}{\linewidth}
\begin{corrige}{6}{Lire un code}
\begin{questions}
\item B C A.
\item Ils indiquent les groupes de deux bits et empêchent une lecture ambiguë ; un protocole peut aussi fixer cette longueur.
\end{questions}
\end{corrige}
\end{minipage}
\par\noindent\begin{minipage}{\linewidth}
\begin{corrige}{7}{Compter des messages}
\begin{questions}
\item Quatre.
\item 00, 01, 10 et 11.
\end{questions}
\end{corrige}
\end{minipage}
\par\noindent\begin{minipage}{\linewidth}
\begin{corrige}{8}{Un étalonnage simple}
\begin{questions}
\item $U=1{,}2\unite V$.
\item $x=0{,}8/0{,}10=8$ dans l’unité de la grandeur donnée.
\end{questions}
\par\smallskip\noindent\textbf{Pour comprendre.}\enspace Pour x = 12, $U=0{,}10\times12=1{,}2\unite V$. Dans l’autre sens, $x=U/0{,}10=0{,}8/0{,}10=8$ dans l’unité de x. Vérifier le domaine d’étalonnage avant d’interpréter.\par
\end{corrige}
\end{minipage}
\par\noindent\begin{minipage}{\linewidth}
\begin{corrige}{9}{Une sortie affine}
\begin{questions}
\item $U=0{,}90\unite V$.
\item Non : à $T=0$, U vaut 0,50 V. C’est une relation affine dans le domaine fourni.
\end{questions}
\par\smallskip\noindent\textbf{Pour comprendre.}\enspace À 20 °C : $U=0{,}020\times20+0{,}50=0{,}90\unite V$. À 0 °C, U vaut déjà 0,50 V ; la droite ne passe donc pas par l’origine : la relation est affine, pas proportionnelle.\par
\pcFigureEx{figures/cor-etalonnage.pdf}{Une droite qui ne passe pas par l’origine n’exprime pas une proportionnalité entre U et T.}
\end{corrige}
\end{minipage}
\par\noindent\begin{minipage}{\linewidth}
\begin{corrige}{10}{Une valeur reçue}
\begin{questions}
\item $T=(1{,}30-0{,}50)/0{,}020=40\,^{\circ}\mathrm C$.
\item Oui, elle appartient à la plage annoncée.
\end{questions}
\par\smallskip\noindent\textbf{Pour comprendre.}\enspace De $U=0{,}020T+0{,}50$, déduire $T=(U-0{,}50)/0{,}020$. Pour 1,30 V, $T=0{,}80/0{,}020=40\,^{\circ}\mathrm C$, dans la plage 0-50 °C.\par
\pcFigureEx{figures/cor-etalonnage-lecture.pdf}{Enlever d’abord la tension de décalage, puis diviser par le coefficient 0,020 V/°C.}
\end{corrige}
\end{minipage}
\par\noindent\begin{minipage}{\linewidth}
\begin{corrige}{11}{Une erreur de bit}
\begin{questions}
\item D.
\item Non : 11 est un code valide. Il faut une information de contrôle supplémentaire pour repérer certains défauts.
\end{questions}
\end{corrige}
\end{minipage}
\par\noindent\begin{minipage}{\linewidth}
\begin{corrige}{12}{Répéter et comparer}
\begin{questions}
\item Au moins une transmission ou lecture est incohérente ; le deuxième groupe diffère.
\item Non à elle seule : la répétition détecte ici un désaccord, mais ne désigne pas la bonne version.
\end{questions}
\end{corrige}
\end{minipage}
\par\noindent\begin{minipage}{\linewidth}
\begin{corrige}{13}{Hors de la plage}
\begin{questions}
\item $T=(3-0{,}50)/0{,}020=125\,^{\circ}\mathrm C$.
\item Ce résultat est hors de la plage validée ; le capteur peut saturer ou changer de comportement. Il faut une autre méthode adaptée.
\end{questions}
\par\smallskip\noindent\textbf{Pour comprendre.}\enspace Le calcul donne $T=(3-0{,}50)/0{,}020=125\,^{\circ}\mathrm C$. Le modèle n’est validé que de 0 à 50 °C : l’extrapolation ne suffit pas à annoncer une température mesurée fiable.\par
\pcFigureEx{figures/cor-hors-plage.pdf}{125 °C est un résultat d’extrapolation, pas une mesure validée par l’étalonnage 0-50 °C.}
\end{corrige}
\end{minipage}
\par\noindent\begin{minipage}{\linewidth}
\begin{corrige}{14}{Une résolution fournie}
\begin{questions}
\item Pas avec cet affichage en pas de 0,5 °C.
\item Non : exactitude, étalonnage, temps de réponse et conditions de contact interviennent aussi.
\end{questions}
\end{corrige}
\end{minipage}
\par\noindent\begin{minipage}{\linewidth}
\begin{corrige}{15}{Analogique ou numérique}
\begin{questions}
\item La tension est modélisée comme analogique ; la sortie chiffrée est discrète et numérique.
\item Non : une résolution limitée fait partie du dispositif. Il faut en tenir compte dans l’interprétation.
\end{questions}
\pcFigureEx{figures/cor-analogique-numerique.pdf}{L’exemple montre des prélèvements et des arrondis prévus par le dispositif, pas des fautes humaines.}
\end{corrige}
\end{minipage}
\par\noindent\begin{minipage}{\linewidth}
\phantomsection\label{cor-defi-20}\pcSolutionDefi{20}{Le protocole du phare}\begin{questions}
\item Deux secondes, puisque chaque message comporte deux bits.
\item Par exemple un allumage plus long de durée convenue, suivi d’une séparation précise, puis les deux bits. Le récepteur doit connaître cette convention.
\item Une même perturbation peut se reproduire, et deux versions différentes ne permettent pas toujours de choisir la bonne. Il faut tester le protocole et discuter ses limites.
\end{questions}
\pcFigureEx{figures/cor-protocole-phare.pdf}{Un début plus long et une séparation convenue distinguent le silence d’un message 00 ; ici, un exemple avec le message 01.}
\end{minipage}
\par\noindent\begin{minipage}{\linewidth}
\phantomsection\label{cor-activite-20}\pcSousTitre{Pistes pour l’investigation}
Une table à deux bits peut coder quatre messages, mais un début et une séparation sont nécessaires pour distinguer 00 d’une absence de transmission. Deux répétitions différentes détectent un désaccord sans indiquer forcément la bonne version. Une erreur transformant un code valide en un autre peut rester invisible ; un protocole doit préciser ses moyens de contrôle et leurs limites.\par\end{minipage}
